\subsection{Lastenheft (Auszug)}
\label{app:Lastenheft}
Es folgt ein Auszug aus dem Lastenheft mit Fokus auf die Anforderungen:

\subsubsection{Funktionale Anforderungen}
\label{sec:FunktionaleAnforderungen}
Die funktionalen Anforderungen beschreiben die konkreten Funktionen, die das System bereitstellen muss.

\subsubsection*{Benutzerverwaltung}
\begin{itemize}
    \item Anmeldung von Besucher*innen mittels Ticketnummer
    \item Verwaltung eines persönlichen Guthabens
    \item Unterscheidung von Benutzerrollen (Besucher*innen, VIP-Besucher*innen, Standpersonal)
\end{itemize}

\subsubsection*{Bestellsystem}
\begin{itemize}
    \item Anzeige verfügbarer Verkaufsstände und Produkte
    \item Darstellung von Produktverfügbarkeiten und geschätzten Wartezeiten
    \item Durchführung von Bestellungen mit Guthabenabbuchung
    \item Generierung eines eindeutigen QR-Codes pro Bestellung
\end{itemize}

\subsubsection*{Bestellverwaltung}
\begin{itemize}
    \item Anzeige aller Bestellungen für Besucher*innen inklusive Status
    \item Verwaltung von Bestellungen durch Standpersonal
    \item Änderung des Bestellstatus (offen, in Zubereitung, abholbereit, ausgegeben)
    \item Priorisierung von Bestellungen für VIP-Besucher*innen
\end{itemize}

\subsubsection*{Abholprozess}
\begin{itemize}
    \item Anzeige des QR-Codes zur Abholung
    \item Verifikation der Bestellung durch Scannen des QR-Codes
    \item Abschluss der Bestellung nach erfolgreicher Ausgabe
\end{itemize}

\subsubsection*{Bestandsverwaltung}
\begin{itemize}
    \item Verwaltung der Produktbestände pro Stand
    \item Automatische Reduzierung des Bestands nach Bestellung oder Ausgabe
    \item Kennzeichnung nicht verfügbarer Produkte
\end{itemize}

\subsubsection*{Stornierung und Sonderfälle}
\begin{itemize}
    \item Stornierung von Bestellungen bei fehlender Verfügbarkeit
    \item Automatische Rückerstattung des Guthabens
    \item Behandlung nicht abgeholter Bestellungen
\end{itemize}

\subsubsection{Nichtfunktionale Anforderungen}
\label{sec:NichtfunktionaleAnforderungen}
Die nichtfunktionalen Anforderungen beschreiben qualitative Eigenschaften des Systems.

\subsubsection*{Benutzbarkeit (Usability)}
\begin{itemize}
    \item Intuitive und übersichtliche Benutzeroberfläche
    \item Schnelle Erlernbarkeit ohne umfangreiche Einarbeitung
    \item Klare Darstellung von Statusinformationen und Rückmeldungen
\end{itemize}

\subsubsection*{Leistungsanforderungen (Performance)}
\begin{itemize}
    \item Schnelle Verarbeitung von Bestellungen in Echtzeit
    \item Geringe Ladezeiten bei Benutzerinteraktionen
    \item Stabiler Betrieb auch bei hoher Anzahl gleichzeitiger Benutzer
\end{itemize}

\subsubsection*{Zuverlässigkeit}
\begin{itemize}
    \item Korrekte Verarbeitung von Bestellungen ohne Datenverlust
    \item Konsistente Speicherung von Bestell- und Bestandsdaten
\end{itemize}

\subsubsection*{Sicherheit}
\begin{itemize}
    \item Schutz vor unbefugtem Zugriff auf Benutzerdaten
    \item Sichere Verarbeitung von Guthaben und Transaktionen
\end{itemize}

\subsubsection*{Wartbarkeit und Erweiterbarkeit}
\begin{itemize}
    \item Modularer Aufbau des Systems
    \item Möglichkeit zur Erweiterung um zusätzliche Funktionen (z. B. weitere Stände)
\end{itemize}